\documentclass[../main.tex]{subfiles}

\begin{document}


Deep learning has become the cornerstone of modern artificial intelligence, driving unprecedented advancements in fields ranging from autonomous transportation and medical diagnostics to natural language understanding and human-computer interaction. These complex neural networks now possess the ability to see, hear, and reason with a proficiency that often meets or exceeds human capabilities. However, the rapid and widespread deployment of these models has outpaced our understanding of their fundamental security limitations. This progress rests on a fragile foundation, vulnerable to a subtle yet powerful class of exploits known as adversarial attacks.

An adversarial attack involves making small, often human-imperceptible perturbations to a model's input data with the express purpose of inducing a specific, confident, and incorrect output. A carefully crafted noisy pattern, invisible to the naked eye, can cause a state-of-the-art image classifier to mistake a stop sign for a speed-limit sign, with potentially catastrophic consequences in an autonomous vehicle. Similarly, minor alterations to an audio command or a sentence's phrasing can deceive speech recognition systems or manipulate the sentiment of powerful language models. These vulnerabilities are not theoretical edge cases; they represent a systemic security risk that threatens the safe and reliable deployment of AI in critical systems.

\section{Problem Statement}
\label{sec:problem_statement}

The central problem this project addresses is the systemic vulnerability of deep learning models across disparate domains---specifically computer vision and natural language processing---to adversarial attacks. While many studies focus on a single task or model, a comprehensive understanding of how these threats manifest, scale, and can be mitigated across different data modalities and model architectures is lacking. This work aims to bridge that gap through a large-scale empirical investigation, providing a comparative analysis of the adversarial threat landscape and the effectiveness of corresponding defenses.

\section{Scope and Objectives}
\label{sec:scope_objectives}

To conduct this investigation, the following specific objectives were established for this project:
\begin{itemize}
    \item To systematically investigate the impact of adversarial attacks on state-of-the-art models for core computer vision tasks, including \textbf{Image Classification} (EfficientNet-B1), \textbf{Image Segmentation} (U-Net), and \textbf{Object Detection} (YOLOv8n).
    \item To extend this investigation to the language domain, evaluating vulnerabilities in \textbf{Large Language Models} (Qwen-1.5), traditional \textbf{Natural Language Processing} for sentiment analysis (BiLSTM-GloVe), and \textbf{Automatic Speech Recognition} systems (Wav2Vec 2.0).
    \item To implement and benchmark a diverse suite of both white-box attacks (e.g., PGD, C\&W, GCG) and black-box attacks (e.g., HopSkipJump, Square Attack) to establish performance degradation baselines across all tasks.
    \item To design, implement, and rigorously evaluate the efficacy of corresponding defense mechanisms, including adversarial training and the development of novel, task-specific defensive pipelines.
    \item To analyze the critical trade-offs between adversarial robustness, model performance on clean data, and the computational latency introduced by each defensive strategy.
\end{itemize}


\section{Key Contributions}
\label{sec:key_contributions}

The primary contributions of this graduation project are:
\begin{itemize}
    \item \textbf{A Broad, Cross-Domain Empirical Analysis:} This work provides a comprehensive and comparative study of adversarial phenomena across six distinct tasks in vision and language, offering a unified baseline for evaluating model robustness.
    \item \textbf{Design of a Novel Dynamic Defensive Pipeline:} For object detection, we designed and implemented a dynamic two-stage defensive pipeline that uses a lightweight adversarial detector to trigger input purification, demonstrating superior robustness against adaptive attacks compared to standard adversarial training.
    \item \textbf{Development of a Hybrid NLP Defense:} For sentiment analysis, we developed a hybrid defense combining Virtual Adversarial Training (VAT) with character-level purification, showing its synergistic effect in mitigating both semantic and character-based attacks.
    \item \textbf{Critical Re-evaluation of ASR Attack Transferability:} Our black-box attack experiments on ASR models provide new evidence that qualifies existing claims in the literature, demonstrating that Self-Supervised Learning (SSL) objective mismatch is a primary barrier to attack transferability.
\end{itemize}

\section{Thesis Structure}
\label{sec:thesis_structure}

The remainder of this report is organized into two main parts. Following a comprehensive literature review of the field in Chapter 2, \textbf{Part I} focuses on Computer Vision. This part contains dedicated chapters on Image Classification (Chapter \ref{chap:classification}), Image Segmentation (Chapter \ref{chap:segmentation}), and Object Detection (Chapter \ref{chap:object_detection}). \textbf{Part II} shifts to Language Models, with chapters covering Large Language Models (LLMs) (Chapter \ref{chap:llm}), Natural Language Processing (NLP) (Chapter \ref{chap:nlp}), and Automatic Speech Recognition (ASR) (Chapter \ref{chap:asr}). Finally, Chapter \ref{chap:conclusion} concludes the report with a summary of our findings, a discussion of their implications, and an outline of promising directions for future research.

\end{document}